% Chapter Template

\chapter{Introduction}\doublespacing % Main chapter title

\label{chap:introduction}

\lhead{Chapter I. \emph{Introduction}}

\setstretch{1.2}
\section{Context and motivation}
In this thesis, we analyze the performance of \gls{cpu} architectures.
Modern \gls{ml} techniques often leverage \gls{ai} models to improve accuracy.

Once introduced, \gls{cpu} will appear only in its short form.

Using classical mechanics, the kinetic energy is defined as shown in Equation~\eqref{eq:kinetic_energy}.

\begin{equation}
    E_k = \frac{1}{2} m v^2
    \label{eq:kinetic_energy}
\end{equation}

\section{Challenges and research questions}
We denote the error tolerance by \ac{eps}. In contrast, the parameter
\ac{avc} is used for controlling the video coding process.

\section{Research methodology}

\subsection{Example Algorithm}
Below is an example of a simple algorithm:

\begin{algorithm}[H]
\caption{Sum of the first $n$ numbers}
\KwIn{Integer $n$}
\KwOut{Sum of integers from $1$ to $n$}
$sum \leftarrow 0$\;
\For{$i \leftarrow 1$ \KwTo $n$}{
    $sum \leftarrow sum + i$\;
}
\Return{$sum$}\;
\end{algorithm}

\subsection{Example Code}
Here is an example of inline code using \texttt{}.

To compute the sum in Python, you can use \texttt{sum(range(1, n+1))}.

For a quick reference on \LaTeX{} math notation, see the cheat sheet:
\url{https://de.wikipedia.org/wiki/Hilfe:TeX}

\section{Contributions}

\section{Outline and roadmap}